%% kapa: exam=2 round=33 year=2022 period=2 subject=appraisal_theory subject_ko=감정평가이론 seq=3 number=3 points=20
%% source: Q-Net 감정평가사 기출문제 (cst00309) artlSeq=5210838 file=raw/1008/5210838/2022년 제33회 감정평가사 2차시험 2교시 문제지(감정평가이론).pdf
%% conversion: kordoc 4.24.51 -> tools/gen2.py (auto); review=proofread-gpt-astra; flags=none
\begin{longquestion}{【문제 3】}
다음 자료를 참고하여 물음에 답하시오. (20점) \par
\begin{framedtext}
<자료> \par
법원감정인인 감정평가사 甲은 손해배상(기) 사건에서 원고가 주장하는 손해액을 구하고 있다. \par
본 사건 부동산(제2종일반주거지역<건폐율 60 \%, 용적률 200 \%>) 매매 당시 매수자인 원고는 부지 내에 차량 2대의 주차가 가능하다는 피고의 주장을 믿고 소유권이전을 완료하였으나, 부지 내의 공간(공지) 부족으로 현실적으로는 주차가 불가능함을 알게 되었다. \par
현장조사 결과 대상 건물(연와조)의 외벽과 인접부동산 담장 사이에 공간이 일부 있으나 협소하여 주차가 불가능한 것으로 나타났다. \par
기준시점 현재 대상 건물은 용적률 110 \%로 신축 후 50년이 경과하였으나 5년 전 단독주택에서 근린생활시설(사무소)로 용도변경 허가를 받은 후 수선을 하여 경제적 잔존내용년수는 10년인 것으로 판단되었다. \par
대상부동산의 인근지역은 기존주택지역에서 소규모 사무실로 변화하는 특성을 보이고 있고 현재 건물의 용도(이용상황)에 비추어 차량 2대의 주차공간 확보가 최유효이용에 해당한다고 조사되었다. \par
\end{framedtext}
\medskip\noindent 물음 1) 이 사안에서 시장자료(market data)를 통하여 손해액을 구하기 위한 감정평가방법과 해당 감정평가방법의 유용성 및 한계점에 대하여 설명하시오. (10점) \par
\medskip\noindent 물음 2) 만일 물음 1)에서 시장자료(market data)를 구할 수 없는 경우, 적용 가능한 다른 감정평가방법들에 대하여 설명하고 이러한 접근방식을 따르는 경우 손해액의 상한은 어떻게 판단하는 것이 합리적인지 설명하시오. (10점) \par
\end{longquestion}
